\documentclass[10pt,twocolumn]{article}
\usepackage[T1]{fontenc}
\usepackage[utf8]{inputenc}
\usepackage{lmodern}
\usepackage[margin=1.8cm,top=2.4cm,bottom=2cm,columnsep=0.6cm]{geometry}
\usepackage{fancyhdr}
\usepackage{titlesec}
\usepackage{booktabs}
\usepackage{amsmath,amssymb,amsfonts}
\usepackage{microtype}
\usepackage{xcolor}
\usepackage{mdframed}
\usepackage{parskip}
\usepackage{enumitem}
\usepackage{hyperref}

\definecolor{rulered}{RGB}{180,30,30}
\definecolor{boxgray}{RGB}{245,245,245}
\definecolor{keywordgray}{RGB}{100,100,100}

\pagestyle{fancy}
\fancyhf{}
\fancyhead[L]{\textsc{\small Claude Mag}}
\fancyhead[C]{\small\textit{Machine Learning Research Digest}}
\fancyhead[R]{\small 2026-09-02}
\fancyfoot[C]{\small\thepage}
\renewcommand{\headrulewidth}{0.8pt}

\titleformat{\section}{\normalsize\bfseries\color{rulered}}{}{0pt}{}%
  [\vspace{-4pt}\color{rulered}\rule{\columnwidth}{0.4pt}]
\titlespacing{\section}{0pt}{8pt}{4pt}

\hypersetup{colorlinks=true,urlcolor=rulered,linkcolor=black}

\begin{document}
\setlength{\parindent}{0pt}
\setlength{\parskip}{4pt}

{\small\color{keywordgray}\textbf{Keywords:}
  LoRA \textbullet{}
  low-rank adaptation \textbullet{}
  Transformer attention \textbullet{}
  approximation theory}\par\vspace{4pt}

{\LARGE\bfseries\raggedright
  How Much Rank Does LoRA Need?
  Rank-Error Bounds for Transformer Attention}\par\vspace{4pt}

{\large\itshape\raggedright
  A Task-Dependent Theory Ties LoRA Rank to
  Downstream-Weighted Tail Energy — and Shows Softmax
  Saturation Can Shrink the Required Rank}\par\vspace{6pt}

{\small\textbf{By} Gerard Conangla Planes}\par\vspace{2pt}

{\small\color{keywordgray}arXiv:2608.26052 \textbullet{} Preprint, August 2026}
\par\vspace{2pt}\noindent{\color{rulered}\rule{\columnwidth}{0.6pt}}\vspace{6pt}

LoRA is the dominant paradigm for adapting large language models to
downstream tasks, yet practitioners select the rank $r$ by heuristic.
This paper provides the first task-dependent theory of LoRA approximation
error for Transformer attention: fixing a pretrained head, a target
attention function, and a downstream input distribution, it bounds the
minimum expected KL divergence achievable by a rank-$r$ query LoRA update
as an explicit function of the downstream-weighted tail energy of the target
update matrix. A striking corollary is that softmax saturation can make the
required rank strictly smaller than the rank needed to match the attention logits.

\begin{mdframed}[backgroundcolor=boxgray,linecolor=rulered,linewidth=1pt,
    innerleftmargin=6pt,innerrightmargin=6pt,
    innertopmargin=6pt,innerbottommargin=6pt]
\textbf{\small AT A GLANCE}\par\vspace{4pt}
\begin{description}[leftmargin=0pt,labelsep=4pt,itemsep=2pt,parsep=0pt,
    font=\small\bfseries]
  \item[Problem:] LoRA rank $r$ is chosen by trial-and-error despite being
    task-dependent; no theory predicts the KL error achievable at rank $r$.
  \item[Why it matters:] Principled rank selection reduces parameter
    overhead while guaranteeing sufficient expressiveness for the adaptation.
  \item[Method:] Bound the minimum expected KL error under rank-$r$ query
    LoRA updates as a function of the downstream-weighted spectral tail energy
    $T_r$ of the target update matrix.
  \item[Result:] Softmax saturation lets rank needed to match the attention
    function be strictly lower than rank needed to match logits — constant
    separation for an explicit family.
  \item[Evidence:] Theoretical bounds, extending to multi-head fused LoRA
    and joint query/key adaptation; construction proving tightness.
\end{description}
\end{mdframed}

\section{The Big Picture}

LoRA (Low-Rank Adaptation) fine-tunes a pretrained weight matrix $W_0$
by adding $\Delta W = BA$ with $B \in \mathbb{R}^{d \times r}$ and
$A \in \mathbb{R}^{r \times k}$, keeping $W_0$ frozen. The rank $r$
controls both the parameter count and the class of reachable adaptations.
In current practice, $r$ is chosen from a small grid $\{1, 2, 4, 8, 16\}$
by cross-validation, with no principled basis for the choice.

For Transformer attention, the query matrix $W_Q$ determines the inner
product structure of the attention mechanism. A rank-$r$ update to $W_Q$
shifts the key-query inner products in at most $r$ directions. If the
downstream task requires shifts in many independent directions, a small $r$
is insufficient; if the task requires shifts in a low-dimensional subspace,
a large $r$ wastes parameters.

This paper asks: given the pretrained $W_Q$, the target $W_Q^*$, and the
downstream input distribution $\mathcal{D}$, what is the minimum KL
divergence between adapted and target attention distributions achievable
by any rank-$r$ query LoRA update?

\section{Why Existing Approaches Fall Short}

Existing analyses of LoRA operate at the level of weight-matrix
approximation: a rank-$r$ update approximates $\Delta W^* = W_Q^* - W_Q$
in spectral or Frobenius norm. But this misses two effects:

First, attention KL divergence depends on \emph{downstream-weighted}
approximation error. Components of $\Delta W^*$ that are large in matrix
norm but project onto low-probability input directions have little effect
on the adapted attention distribution. A matrix-norm-based analysis
overestimates the required rank for these components.

Second, softmax applies a highly nonlinear map to the logits before
computing the attention weights. Large logit differences saturate the
softmax and make the attention distribution insensitive to further
increases. A low-rank logit update can therefore achieve the same
attention KL as a higher-rank update that matches the full logit change,
because the relevant logit differences are already in the saturation regime.

\section{Core Mechanism}

The paper defines the \textbf{tail energy} $T_r$ as the sum of eigenvalues
beyond rank $r$ in the downstream-weighted spectral decomposition of the
target update:
\[
  T_r = \sum_{j > r} \lambda_j\!\left(\Sigma^{1/2} \Delta W^{*\top}
        \Delta W^* \Sigma^{1/2}\right)
\]
where $\Sigma = \mathbb{E}_{\mathcal{D}}[x x^\top]$ is the downstream
input covariance. The tail energy measures the fraction of the target
adaptation that lies outside the subspace reachable by rank-$r$ queries
under the downstream distribution.

This quantity bridges weight-matrix approximation and attention KL
divergence: a small $T_r$ implies there exists a rank-$r$ update with
small expected KL error; a large $T_r$ implies no rank-$r$ update can
achieve small error.

\section{Technical Formulation}

Fix an input $x$, query $q = W_Q x$, target query $q^* = W_Q^* x$, and
key vectors $\{k_j\}$. The logit residual is $d = q - q^*$.

The KL divergence between the adapted and target attention distributions
is:
\[
  \text{KL}(\text{softmax}(q^*) \| \text{softmax}(q^* - \Delta q))
  = \psi(\|d\|_2)
\]
where $\psi$ is an increasing function bounded above by
$\min\{\|d\|_2^2/4, \sqrt{2}\|d\|_2\}$.

For attention probabilities bounded away from zero (condition $\pi_j \geq
\gamma > 0$ for all $j$), the lower bound is:
\[
  \text{KL} \geq \psi(\|d\|_2) \geq c_\gamma \|d\|_2
\]
for an explicit constant $c_\gamma$ depending on $\gamma$.

Under realizability, geometry, and moment conditions, the best rank-$r$
expected KL error $\epsilon^*(r)$ satisfies:
\[
  c_1 \psi(\sqrt{T_r}) \;\leq\; \epsilon^*(r) \;\leq\;
  \min\!\left\{\tfrac{T_r}{4},\, \sqrt{2 T_r}\right\}
\]
with explicit constants $c_1$. These bounds are tight: the paper constructs
families achieving both.

\section{Learning or Inference Procedure}

The theory yields a practical rank selection procedure:
\begin{enumerate}[itemsep=2pt,parsep=0pt]
  \item Collect a sample $\{x_i\}$ from the downstream distribution
  \item Estimate $\Sigma = \frac{1}{n}\sum_i x_i x_i^\top$
  \item Compute the downstream-weighted SVD of the target update
    $\Delta W^*$ (estimated from a small fine-tuning probe if needed)
  \item Evaluate $T_r$ for each candidate rank $r$
  \item Select $r$ as the smallest rank with $T_r$ below the
    desired KL error tolerance
\end{enumerate}

At fine-tuning time, only the rank-$r$ query LoRA is trained; all other
weights remain frozen.

\section{What the Guarantee Says}

The bounds establish that the minimum KL error is controlled by $T_r$:
if $T_r \leq \tau^2$, then a rank-$r$ update can match the target
attention to within $\min\{\tau^2/4, \sqrt{2}\tau\}$ expected KL. For
saturated attention (large logit differences), the softmax nonlinearity
allows this bound to be substantially tighter than the logit approximation
error $T_r$ alone would suggest.

The paper proves that there exist target functions where the attention-matching
rank is a constant factor smaller than the logit-matching rank—closing a
gap between empirical observations that low rank often suffices and the
prior absence of theory explaining why.

\section{Experimental Findings}

The paper validates the rank-error bounds on two settings:

\textbf{Synthetic attention heads:} For planted low-rank targets, the
measured KL error at each rank closely tracks $\min\{T_r/4, \sqrt{2T_r}\}$,
confirming the upper bound is tight within a factor of 2.

\textbf{LLM fine-tuning (RoBERTa-large, GLUE):} The rank selected by the
tail energy criterion achieves comparable task accuracy to a grid-searched
rank in 8/9 GLUE tasks, while using 30--60\% fewer LoRA parameters on
average. On STS-B (a semantic similarity task where attention must capture
fine-grained input directions), the criterion correctly predicts that rank
4 is insufficient and recommends rank 16.

\section{Ablations and Interpretation}

\textbf{Softmax saturation effect:} For heads with large pre-softmax
logit magnitudes ($\|q^*\|_\infty > 10$), the attention-matching rank is
consistently 2--4$\times$ lower than the logit-matching rank, validating
the theoretical separation.

\textbf{Multi-head fused LoRA:} When a single shared $B$ is used across
all heads, $T_r$ must aggregate over head-specific downstream covariances;
the bound remains tight under this extension.

\textbf{Joint query/key updates:} The analysis generalizes to adapting
both $W_Q$ and $W_K$ simultaneously, with $T_r$ defined over the joint
key-query interaction tensor.

\textbf{Distribution shift:} If $\mathcal{D}$ is misspecified (e.g.,
using a generic corpus when the downstream distribution is specialized),
$T_r$ may underestimate the required rank. Domain-specific calibration
data is therefore recommended for rank estimation.

\section*{Reference}

Gerard Conangla Planes.
\textbf{How Much Rank Does LoRA Need? Rank-Error Bounds for
Transformer Attention.}
arXiv:2608.26052, August 2026.
\url{https://arxiv.org/abs/2608.26052}

\end{document}
